%!TEX program = xelatex
%!TEX enableShellEscape = true
%!TEX format = pdf
%!TeX spellcheck = en-GB
%!TEX root = introduction.tex
% -- Author: Jannes Bantje, j.bantje@wwu.de
\documentclass[a4paper,index=totoc,toc=bibliography,fontsize=11pt,DIV=13,headinclude,BCOR=5mm,cleardoublepage=empty,draft]{scrartcl}
\usepackage{scrpage2} % wie fancyhdr, nur optimiert auf KOMA-Skript, leicht andere Syntax
\usepackage{etoolbox,letltxmacro} % zum Programmieren
\usepackage[final]{graphicx}
\robustify\rotatebox

% -- Farben
\usepackage[x11names]{xcolor}
\definecolor{dark_gray}{gray}{0.45}
\definecolor{light_gray}{gray}{0.6}
\definecolor{fb10_blue}{cmyk}{0.8,0.4,0.13,0.07}

\usepackage{mathtools} % beinhaltet amsmath
\mathtoolsset{centercolon,showmanualtags}
\newtagform{brackets}{[}{]}
\usetagform{brackets}
\usepackage{fix-cm}
% \usepackage{wasysym} % zusätzliche Symbole
\usepackage[bbgreekl]{mathbbol}
\usepackage{amssymb} % zusätzliche Symbole
% \usepackage{stmaryrd} % für Blitz
\usepackage{nicefrac} % schräge Brüche
\usepackage{faktor}
\newcommand{\Faktor}[1]{\faktor[\textstyle]{#1}}
\usepackage{xfrac}
\usepackage{cancel}
\usepackage{mathdots} % Verbesserung von Punkten wie zB \ldots
\usepackage[bb=px]{mathalfa}
\usepackage{centernot}

\usepackage{silence}
\WarningFilter{latexfont}{Size substitutions with differences}
\WarningFilter{latexfont}{Font shape `U/bbold/m/n' in size}

%!TEX root = introduction.tex
\DeclareMathOperator{\sgn}{sgn} % Signum
\DeclareMathOperator{\im}{im} % Bild
\DeclareMathOperator{\supp}{supp} % Träger/Support
\DeclareMathOperator{\id}{id} % Identität
\DeclareMathOperator{\pr}{pr} % Projektion
\DeclareMathOperator{\Hom}{Hom} % Homomorphismen
\DeclareMathOperator{\sym}{sym} % symmetrisch
\DeclareMathOperator{\Skew}{skew} % schiefsymmetrisch
\DeclareMathOperator{\coker}{coker}
% \DeclareMathOperator{\deg}{deg}
% \DeclareMathOperator{\zentrum}{Z}
\newcommand{\zentrum}{\mathrm{Z}}

% -- Matrixgruppen
\DeclareMathOperator{\Mat}{Mat} % Matrizen
\DeclareMathOperator{\End}{End} % Endomorphismen
\DeclareMathOperator{\GL}{GL} % allgemeine lineare Gruppe
\DeclareMathOperator{\SL}{SL} % spezielle lineare Gruppe
\DeclareMathOperator{\On}{O} % orthogonale Gruppe
\DeclareMathOperator{\SO}{SO} % spezielle orthogonale Gruppe
\DeclareMathOperator{\Un}{U} % unitäre Gruppe
\DeclareMathOperator{\SU}{SU} % spezielle unitäre Gruppe

% -- besserer underbrace befehl
\newcommand{\Underbrace}[2]{{\underbrace{#1}_{#2}}}


% -- Zum Finetuning von Befehlen
\makeatletter
\newcommand{\raisemath}[1]{\mathpalette{\raisem@th{#1}}}
\newcommand{\raisem@th}[3]{\raisebox{#1}{$#2#3$}}
\makeatother


%--Mengen
\newcommand\SetSymbol[1][]{\nonscript\:#1\vert\allowbreak\nonscript\:\mathopen{}}
\providecommand\given{} % to make it exist
\DeclarePairedDelimiterX\set[1]\{\}{\renewcommand\given{\SetSymbol[\delimsize]}#1}

% -- Betrag und Skalarprodukt
\DeclarePairedDelimiter{\abs}{\lvert}{\rvert}
\DeclarePairedDelimiterX\skal[2]{\langle}{\rangle}{#1\,,\,#2}

%--Umklammern
\DeclarePairedDelimiter\enbrace{(}{)}
\DeclarePairedDelimiter\benbrace{[}{]}
\DeclarePairedDelimiter\homo{\llbracket}{\rrbracket}
\newcommand{\ssbrace}[1]{{\scriptscriptstyle\enbrace{#1}}}

%--Norm
\DeclarePairedDelimiter\norm{\Vert}{\Vert}

% -- Box mit vorgegebener minimaler Länge
\DeclareRobustCommand{\minwidthbox}[2]{%
  \ifmmode
    \expandafter\mathmakebox
  \else
    \expandafter\makebox
  \fi
  [\ifdim#2<\width\width\else#2\fi]{#1}%
}

% -- farbiges Untersteichen im Mathe-Modus
\def\mathul#1#2{\color{#1}\underline{{\color{black}#2}}\color{black}} 


% - verbessertes Integral
\newcommand{\Int}[1]{\int_{\mathrlap{#1}}\,\,}

% -- Alles spezielles aus der Differentialtopologie
\newcommand{\Tang}{\ensuremath{\mathrm{T}\mkern-0.85mu}}
\newcommand{\mathd}{\ensuremath{\mathrm{d}\mkern-0.7mu}}
\newcommand{\diff}[2]{\ensuremath{\frac{{\partial #1}}{{\partial #2}}}}
\newcommand{\diffs}[2]{\ensuremath{\partial #1/\partial #2}}
\newcommand{\diffd}[2]{\ensuremath{\frac{\mathd #1}{\mathd #2} }}
\DeclareMathOperator{\Lie}{L}
\DeclareMathOperator{\Ad}{Ad}  
\DeclareMathOperator{\ad}{ad}
\DeclareMathOperator{\Spin}{Spin}
% \DeclareMathOperator{\spin}{spin}
\newcommand{\spin}{\mathop{\mathfrak{spin}}}
\newcommand{\Ce}{\mathcal{C}}
\DeclareMathOperator{\vol}{vol}


% -- Kategorien und Ähnliches
\newcommand{\VEKT}{\textsc{Vekt}}

%--Abbildungsdefinition
\newcommand{\mapdef}[5]{%
	\[
		\begin{array}{rcl}
			\textstyle #1 &\xrightarrow{\minwidthbox{#5}{2em}} & \textstyle #2 \\[0.5ex]
			\textstyle #3 &\xmapsto{\minwidthbox{\mbox{ }}{2em}} & \textstyle #4
		\end{array}
	\]
}

%--modifiziertes Stackrel
\newcommand{\StackText}[2]{\stackrel{\mbox{\scriptsize #1}}{#2}}
\newcommand{\StackTextClap}[2]{\stackrel{\mathclap{\mbox{\scriptsize #1}}}{#2}}
 % Liste mit Mathebefehlen laden
% \allowdisplaybreaks

% -- TikZ-Kram
\usepackage{tikz-cd}
\usetikzlibrary{quotes,babel,arrows.meta}
\tikzset{>=Straight Barb[]}
\tikzcdset{arrow style=math font}

% -- tikz/externalize
\usetikzlibrary{external}
\tikzexternalize[prefix=Bilder/tikz/, up to date check=diff]
\tikzset{external/system call={xelatex \tikzexternalcheckshellescape %-- verwende LuaLaTeX, wegen dynamischer Speicherallokation
    -halt-on-error -interaction=batchmode -jobname "\image" "\texsource"}}
\AtBeginEnvironment{tikzcd}{\tikzexternaldisable}
\AtEndEnvironment{tikzcd}{\tikzexternalenable}
\pgfkeys{/pgf/images/include external/.code=\includegraphics{#1}}

% -- Einstellungen bezüglich Schriftart etc.
% \usepackage{newpxtext}
\usepackage{palatino}
\usepackage[bigdelims,varbb,varg]{newpxmath}
% \usepackage[euler-digits,euler-hat-accent]{eulervm} % Mathe-Schriftart
\usepackage[no-math]{fontspec}
\defaultfontfeatures{Mapping=tex-text, Ligatures={Required,Common,Contextual}} % Word Spacing?

% \setmainfont{Syntax Next LT W1G}[UprightFont=* Regular, BoldFont=* Bold, BoldItalicFont=* Bold Italic, ItalicFont= * Italic]
% \setmainfont{texgyrepagella}[Extension=.otf,UprightFont=*-regular,BoldFont=*-bold,ItalicFont=*-italic,BoldItalicFont=*-bolditalic]
% \setmainfont{Optima}
% \setmainfont{Source Sans Pro}[Ligatures=Common, BoldFont={* Bold}, ItalicFont={* Light Italic}] % ,SmallCapsFont={Linux Libertine O}, SmallCapsFeatures={Letters=SmallCaps}
\setmainfont{TeXGyrePagellaX}[Extension=.otf,UprightFont=*-Regular,BoldFont=*-Bold,ItalicFont=*-Italic,BoldItalicFont=*-BoldItalic,ItalicFeatures={Style=Historic},Ligatures={Required,Common,Contextual,Historic},Numbers=OldStyle]

% \setsansfont{OfficinaSerifBookC}[ItalicFont=*-Italic,BoldFont=OfficinaSerifBoldC,BoldItalicFont=OfficinaSerifBoldC-Italic]
% \setsansfont{DIN Next Slab Pro}[Scale=MatchLowercase,BoldFont=* Medium, BoldItalicFont=* Medium Italic]
% \setsansfont{Syntax Next LT W1G}[UprightFont=* Regular, BoldFont=* Bold, BoldItalicFont=* Bold Italic, ItalicFont= * Italic]
% \setsansfont{texgyreadventor}[Scale=MatchUppercase, Extension=.otf, UprightFont=*-regular, BoldFont=*-bold, ItalicFont=*-italic,BoldItalicFont=*-bolditalic]
% \setsansfont{texgyrepagella}[Scale=MatchUppercase,Extension=.otf,UprightFont=*-regular,BoldFont=*-bold,ItalicFont=*-italic,BoldItalicFont=*-bolditalic]
\setsansfont{DIN Next W1G}[UprightFont=* Regular, BoldFont=* Bold, BoldItalicFont=* Bold Italic, ItalicFont= * Italic]

% \setmonofont{Monaco}
\setmonofont{SourceCodePro}[Scale=MatchLowercase,Extension=.otf,UprightFont=*-Regular, BoldFont=*-Semibold, ItalicFont=*-Light]
% \setmonofont{texgyrecursor}[Scale=MatchLowercase, Extension=.otf, UprightFont=*-regular, BoldFont=*-bold, ItalicFont=*-italic,BoldItalicFont=*-bolditalic]
% \newfontface\fancy{Insignia LT Std Roman}[Scale=MatchUppercase,BoldFont=*,ItalicFont=*,BoldItalicFont=*,] % Alternate=1
\addtokomafont{disposition}{\sffamily}
\linespread{1.07}
% \addtokomafont{chapter}{\fancy}
% \addtokomafont{chapterentry}{\fancy}
% \addtokomafont{part}{\fancy}
% \addtokomafont{partentry}{\fancy}
% \addtokomafont{section}{\fancy}
% \addtokomafont{title}{\fancy\addfontfeature{Kerning=Uppercase,LetterSpace=8}}
\defaultfontfeatures{Mapping=tex-text, Ligatures={Required,Common,Contextual},Extension=.otf}
\usepackage{fontawesome}
\usepackage[final]{microtype}
\flushbottom
\KOMAoptions{DIV=last}


% -- Einstellungen bezüglich der Sprache, Silbentrennung etc.
\usepackage{polyglossia}
\setmainlanguage[variant=british]{english}

% -- BibLaTeX
\usepackage[%
	backend=biber,
	sortlocale=auto,
	natbib,
	hyperref,
	backref,
	style=alphabetic
	]%
{biblatex}
\renewcommand*{\mkbibnamefamily}[1]{%
  \ifmknamesc{\textsc{#1}}{#1}}

\renewcommand*{\mkbibnameprefix}[1]{%
  \ifboolexpr{ test {\ifmknamesc} and test {\ifuseprefix} }
    {\textsc{#1}}
    {#1}}

\def\ifmknamesc{%
  \ifboolexpr{ test {\ifcurrentname{labelname}}
               or test {\ifcurrentname{author}}
               or ( test {\ifnameundef{author}} and test {\ifcurrentname{editor}} ) }}
\addbibresource{quellen.bib}

% -- Konfiguration von Hyperref
\usepackage[hidelinks, pdfpagelabels, bookmarksopen=true, bookmarksnumbered=true, linkcolor=black, urlcolor=SkyBlue2, plainpages=false,pagebackref, citecolor=black, hypertexnames=true, pdfauthor={Jannes Bantje}, pdfborderstyle={/S/U}, linkbordercolor=SkyBlue2, colorlinks=false,backref=false]{hyperref}
\hypersetup{final}

\usepackage[nameinlink,noabbrev]{cleveref}

% -- \url aufhübschen
\newcommand{\appendLink}[1]{#1\,\faExternalLink}
\newcommand{\hrefsym}[2]{\href{#1}{\texttt{\appendLink{#2}}}}
\renewcommand{\url}[1]{\hrefsym{#1}{\nolinkurl{#1}}}

% -- Aufzählungen, Anführungszeichen etc.
\usepackage[shortlabels,inline]{enumitem}
\setlist[itemize,1]{label=\faCaretRight}
\setlist[enumerate]{font=\bfseries}
\setlist{itemsep=0pt}
\usepackage[autostyle,german=quotes]{csquotes}


% -- Indexverarbeitung
\usepackage{makeidx}
\newcommand{\bet}[1]{\textbf{\emph{#1}}}
\newcommand{\Index}[1]{\bet{#1}\index{#1}}
\makeindex
% \setindexpreamble{{\noindent \itshape Die Seitenzahlen sind mit Hyperlinks zu den entsprechenden Seiten versehen, also anklickbar}\par \bigskip}
\renewcommand{\indexpagestyle}{scrheadings}

% -- Alles zum Thema Gliederung
% \usepackage{chngcntr}
% \counterwithin*{equation}{section}
\numberwithin{equation}{section}
% \setcounter{tocdepth}{3}

% -- theorem packages
\let\openbox\relax
\usepackage{amsthm}
\usepackage{thmtools,thm-restate}
\usepackage{mdframed}
\renewcommand{\listtheoremname}{Übersicht aller Aussagen}

% -- Theoreme als PDF-Lesezeichen
\usepackage{bookmark}
\bookmarksetup{open,numbered}
\makeatletter
\newcommand*{\theorembookmark}{%
  \bookmark[
    dest=\@currentHref,
    rellevel=1,
    keeplevel,
  ]{%
    \thmt@thmname\space\csname the\thmt@envname\endcsname
    \ifx\thmt@shortoptarg\@empty
    \else
      \space(\thmt@shortoptarg)%
    \fi
  }%
}
\makeatother

% -- Definition der einzelnen Umgebungen
\declaretheoremstyle[%
	headfont=\sffamily\bfseries,
	notefont=\normalfont\sffamily,
	bodyfont=\normalfont,
	headformat=\NUMBER\ \NAME\NOTE,
	headpunct=.,
	postheadspace=1em,
	spaceabove=15pt,spacebelow=10pt,
	shaded={bgcolor=gray!20},
	% mdframed={%
	% 	backgroundcolor=gray!20,
	% 	innertopmargin=0pt,
	% 	innerleftmargin=0pt,
	% 	% roundcorner=5pt,
	% 	innerbottommargin=0pt,
	% 	splittopskip = 0pt,
	% 	% skipbelow=6pt,
	% 	% skipbelow=6pt,
	% 	topline=false,bottomline=false,leftline=false,rightline=false
	% },
	postheadhook=\theorembookmark]%
{mainstyle}
\declaretheoremstyle[%
	headfont=\sffamily\bfseries,
	notefont=\normalfont\sffamily,
	bodyfont=\normalfont,
	headformat=\NUMBER\ \NAME\NOTE,
	headpunct=.,
	postheadspace=1em,
	spaceabove=15pt,spacebelow=10pt,
	shaded={bgcolor=fb10_blue!20},
	postheadhook=\theorembookmark]%
{mainstyle_blue}
\declaretheoremstyle[%
	headfont=\sffamily\bfseries,
	notefont=\normalfont\sffamily,
	bodyfont=\normalfont,
	headformat=\NUMBER\ \NAME\NOTE,
	headpunct=.,
	postheadspace=1em,
	spaceabove=15pt,spacebelow=10pt,
	postheadhook=\theorembookmark]%
{mainstyle_unshaded}
\declaretheoremstyle[%
	headfont=\sffamily\bfseries,
	notefont=\normalfont\sffamily,
	bodyfont=\normalfont,
	headformat=\NUMBER\NAME\NOTE,
	headpunct=.,
	postheadspace=1em,
	spaceabove=15pt,spacebelow=10pt,
	% shaded={bgcolor=gray!20},
	postheadhook=\theorembookmark]%
{mainstyle_unnumbered}
\declaretheoremstyle[%
	headfont=\sffamily\bfseries,
	notefont=\normalfont\sffamily,
	bodyfont=\normalfont,
	headformat=swapnumber,
	headpunct=.,
	postheadspace=1em,
	spaceabove=15pt,spacebelow=10pt,
	shaded={bgcolor=gray!20},
	postheadhook=\theorembookmark,
	qed=\qedsymbol]%
{mainstyleB}
\declaretheorem[name=Definition,parent=section,style=mainstyle_blue]{definition}
\declaretheorem[name=Definition \& Proposition,refname=Proposition,sharenumber=definition,style=mainstyle_blue]{definitionP}
\declaretheorem[name=Definition,numbered=no,style=mainstyle_unnumbered]{definition*}
\declaretheorem[name=Theorem,sharenumber=definition,style=mainstyle]{theorem}
\declaretheorem[name=Theorem,numbered=no,style=mainstyle_unnumbered]{theorem*}
\declaretheorem[name=Proposition,sharenumber=definition,style=mainstyle,refname=Proposition]{proposition}
\declaretheorem[name=Lemma,sharenumber=definition,style=mainstyle]{lemma}
\declaretheorem[name=Satz,refname=Satz,sharenumber=definition,style=mainstyle]{satz}
\declaretheorem[name=Satz,sharenumber=definition,style=mainstyle_unshaded]{satzUnshaded}
\declaretheorem[name=Definition,sharenumber=definition,style=mainstyle_unshaded]{definitionUnshaded}
\declaretheorem[name=Satz,numbered=no,style=mainstyle_unnumbered]{satz*}
\declaretheorem[name=Korollar,sharenumber=definition,style=mainstyle]{korollar}
\declaretheorem[name=Korollar,sharenumber=definition,style=mainstyleB]{korollarB}
\declaretheorem[name=Frage,numbered=no,style=mainstyle_unnumbered]{frage}
\declaretheorem[name=Frage,sharenumber=definition,style=mainstyle_unshaded]{frageA}
\declaretheorem[name=Erinnerung,sharenumber=definition,style=mainstyle_unshaded]{erinnerungA}
\declaretheorem[name=Ausblick,sharenumber=definition,style=mainstyle_unshaded]{ausblick}
\declaretheorem[name=Konvention,sharenumber=definition,style=mainstyle]{konvention}
\declaretheorem[name=Notation,sharenumber=definition,style=mainstyle_unshaded]{notation}
\declaretheorem[name=Bemerkung,sharenumber=definition,style=mainstyle_unshaded]{bemerkung}
\declaretheorem[name=Bemerkung,numbered=no,style=mainstyle_unnumbered]{bemerkung*}
\declaretheorem[name=Beispiel,sharenumber=definition,style=mainstyle_unshaded]{beispiel}
\declaretheorem[name=Beispiel,numbered=no,style=mainstyle_unnumbered]{beispiel*}
\declaretheorem[name=Exkurs,numbered=no,style=mainstyle_unnumbered]{exkurs*}

% english versions
\declaretheorem[name=Remark,sharenumber=definition,style=mainstyle_unshaded]{remark}
\declaretheorem[name=Remark,numbered=no,style=mainstyle_unnumbered]{remark*}
\declaretheorem[name=Example,sharenumber=definition,style=mainstyle_unshaded]{example}
\declaretheorem[name=Corollary,sharenumber=definition,style=mainstyle]{corollary}

% -- Beweise
\declaretheoremstyle[headfont=\bfseries\scshape,bodyfont=\normalfont,headpunct=:,postheadspace=1em,spacebelow=12pt,spaceabove=2pt,qed=\qedsymbol]{beweise}
\declaretheoremstyle[headfont=\bfseries\scshape,bodyfont=\normalfont,headpunct=:,postheadspace=1em,spacebelow=12pt,spaceabove=2pt]{beweisskizze}
\declaretheoremstyle[headfont=\sffamily\bfseries,bodyfont=\normalfont,headpunct=:,postheadspace=1em,spacebelow=10pt,spaceabove=10pt]{bemerkungen}
\declaretheorem[name=Beweis,numbered=no,style=beweise]{beweis}
\let\proof\relax
\declaretheorem[name=Proof,numbered=no,style=beweise]{proof}
\declaretheorem[name=Sketch of Proof,numbered=no,style=beweisskizze]{sketch}

\declaretheorem[name=Übung,numbered=no,style=bemerkungen]{uebung}
\declaretheorem[name=Erinnerung,numbered=no,style=bemerkungen]{erinnerung}


% \makeatletter
% \long\def\theequation{\@fnsymbol \c@equation}
% \makeatother

% -- marginnotes package
\usepackage[fulladjust]{marginnote}
\renewcommand*{\marginfont}{\itshape \footnotesize}
\usepackage{ragged2e}
\renewcommand*{\raggedleftmarginnote}{\RaggedLeft}
\renewcommand*{\raggedrightmarginnote}{\RaggedRight}

% -- todonotes package
\usepackage[textsize=small,obeyDraft]{todonotes}
\LetLtxMacro{\oldtodo}{\todo}
\renewcommand{\todo}[2][]{\tikzexternaldisable\oldtodo[#1]{#2}\tikzexternalenable}
\LetLtxMacro{\oldmissingfigure}{\missingfigure}
\renewcommand{\missingfigure}[2][]{\tikzexternaldisable\oldmissingfigure[{#1}]{#2}\tikzexternalenable}

% -- Fußnoten anpassen
\deffootnote[1.5em]{1.5em}{1.5em}{\textsuperscript{\thefootnotemark}\ }

% -- Figures und Captions
\usepackage{wrapfig}

%--Römische Zahlen
\newcommand{\RM}[1]{\MakeUppercase{\expandafter{\romannumeral #1.\relax}}}
\robustify\RM


% -- Metadaten für die Titelei
\author{Jannes Bantje}


% -- Kopf- und Fußzeilen
\setheadsepline{1pt}[\color{light_gray}]
\pagestyle{scrheadings}
\clearscrheadfoot
\automark{section}
\lehead{\includegraphics[height=0.6 cm,keepaspectratio]{Bilder/Logo_WWU_Muenster_light_gray.pdf}}
\lohead{\color{dark_gray}\normalfont\footnotesize\sffamily\itshape Introduction to Knots and Primes -- Jannes Bantje}
\rohead{\includegraphics[height=0.6 cm,keepaspectratio]{Bilder/fb10logo_gray.pdf}}

\ofoot[{\color{dark_gray}\sffamily\LARGE \thepage}]{{\color{dark_gray}\sffamily\LARGE \thepage}} %hier wir auch der plain Stil bearbeitet!
\ifoot{ \color{dark_gray} \sffamily\small \leftmark}

% -- Inhaltsverzeichnis und weitere Pakte, die zuletzt geladen werden sollten
\usepackage[tocindentauto]{tocstyle}
\usetocstyle{KOMAlike}
\usepackage{ellipsis}

% \addbibresource{quellen.bib}

\begin{document}
\pagenumbering{Roman}

%!TEX root = introduction.tex
\begin{titlepage}
\hspace*{0.15\textwidth} % Whitespace to the left of the title page
\rule{2pt}{\textheight} % Vertical line
\hspace*{0.05\textwidth} % Whitespace between the vertical line and title page text
\begin{minipage}[b]{0.75\textwidth}
	\raggedright
	\includegraphics[height=1.5cm, keepaspectratio]{Bilder/Logo_WWU_Muenster.pdf} \\[2cm]
	{\Large \sffamily Knots, Primes and other Analogies}\\[0.5cm]
	{\Huge\sffamily\bfseries Introduction to\\[1ex] Knots and Primes}\\[7cm]
	{\large \textit{Author:}}\\[5pt]
	{\Large \textsc{Jannes Bantje}}\\[5pt] % Author name
	{\small\nolinkurl{j.bantje@wwu.de}}\\[2.5cm]
	
	\vspace{0.1\textheight}
	\includegraphics[height=1cm, keepaspectratio]{Bilder/fb10logo.pdf}
	% \hfill\includegraphics[height=8cm,keepaspectratio]{Bilder/E8_graph.pdf}
\end{minipage}	


% \begin{center}
% \vspace*{10cm}
% \includegraphics[height=9cm,keepaspectratio]{Bilder/E8_graph.pdf}
% \end{center}

\end{titlepage}


\tableofcontents

\tableofcontents
\section*{Notation}
\addcontentsline{toc}{section}{Notation}

\cleardoubleoddemptypage%

\pagenumbering{arabic}
\setcounter{page}{1}
\setcounter{footnote}{0}

\section{Introduction to knot theory} % (fold)
\label{sec:intro_knots}

\begin{definition}
	A \Index{knot} is the image of an embedding $S^1 \hookrightarrow S^3$.
	A \Index{knot type} is the equivalence class of such embeddings under ambient isotopy.
\end{definition}

We usually wont distinguish between these two notions and refer to a knot type simply as a knot.
To exclude some pathological knots we restrict our attention to \emph{tame knots}:

\begin{definition}
	A knot $K$ is called \bet{tame}\index{knot!tame} if it is ambient isotopic to a polygonal knot, i.e. a knot that is piecewise linear.
\end{definition}

\begin{example}
	Every smooth knot is tame, because by a classical theorem of calculus every $C^1$ map $\gamma \colon [0,1] \to \mathbb{R}^n$ is rectifiable and this can be used to construct an ambient isotopy, see \textcite[Appendix I.]{introduction-knots} (very technical!).
	It is also possible to show, that a polygonal knot can be transformed into a smooth knot by an ambient isotopy.
\end{example}

\textcite[p.~9]{morishita} assumes every topological space considered to be a PL-manifold and every map between topological spaces to be PL as well.

Since we are interested only in tame knots, for every knot $K$ we have a \Index{tubular neighbourhood} $V_k$, that is: There is a vector bundle $E \to K$ and an embedding $f \colon E \to S^3$ such that (with $s_0$ being the zero section of $E \to K$)
\begin{itemize}
	\item $f \circ s_0 = {\id_K}$ and
	\item $V_K \coloneqq f(E)$ is an open neighbourhood of $K$ in $S^3$. \hfill see \textcite{hirschDifferential}
\end{itemize}
Note that any two such tubular neighbourhoods are isotopic, see \textcite[Chap.~4, Thm.~5.3]{hirschDifferential}.

\begin{definition}
	Let $K$ be a knot. $X_K \coloneqq S^3 \setminus \interior(V_k)$ is the \Index{knot exterior} of $K$.
	The \Index{knot group} is defined as the fundamental group of the knot exterior: $G_K \coloneqq \pi_1(X_K)$.
\end{definition}

\begin{lemma}
	The knot group is a \Index{knot invariant}, i.e. for two ambient isotopic Knots $K_1$ and $K_2$ we have
	\[
		G_{K_1} \cong G_{K_2}
	\]
\end{lemma}
\begin{proof}
	Pick embeddings $g_1, g_2 \colon S^1 \hookrightarrow S^3$ with $g_i(S^1)=K_i$.
	Let $F \colon S^3 \times I \to S^3$ be an ambient isotopy, i.e. $F_0 = {\id_M}$, $F_t$ is a homeomorphism for each $t \in I$ and $F_1 \circ g_1 = g_2$.
	Now $F_1$ restricted to the knot exterior $X_{K_1}$ is a homeomorphism $X_{K_1} \cong X_{K_2}$ and therefore $G_{K_1} \cong G_{K_2}$ since the fundamental group is a topological invariant.
\end{proof}

\begin{remark}
	\begin{itemize}
		\item The boundary of $X_K$ is a 2-dimensional torus
		\item A \Index{meridian} of $K$ is a closed (oriented) curve which is the boundary of a disk $D^2$ in $V_K$.
		\item A \Index{longitude} is a closed curve in $\partial X_K$ which intersetcts which intersects with a meridian at one point and is null-homologous in $X_K$.
	\end{itemize}
\end{remark}

By removing a point $\infty \in S^3$ we may assume, that a knot $K$ is contained in $\mathbb{R}^3$.
When dealing with knots one is often interested in two-dimensional representations of knots, so called \Index{knot diagrams}:

\begin{definition}
	A projection of $K$ onto a plane in $\mathbb{R}^3$ is called \Index{regular}, if it has only a finite number of multiple points, all of which are double points.
\end{definition}

It's possible to turn any projection of a knot into a regular one by slight pertubation of the knot, so we only consider regular knot projections.
If one marks the overcrossing line in the projection the knot itself can be recovered from the projection.

\begin{definition}[{name=[crossing number]}]
	The \Index{crossing number} of a knot (type) is the least number of crossings in any projection of a knot of that type.
\end{definition}

\begin{definition}
	Given two knots $J$ and $K$ the \Index{connected sum} or \Index{composition} $J\#K$ is the knot obtained by removing a small arc from each knot projection and connecting the endpoint by two new arcs.

	A knot is called \Index{prime}, if it cannot be written as the connected sum of two non-trivial knots and \Index{composite} otherwise.
\end{definition}

\todo[inline]{draw trefoil and other prime knots}

\begin{definition}
	A knot is called alternating if it has a projection in which the crossings alternate between over- and undercrossings while travelling along the knot.
\end{definition}

\begin{lemma}
	Two knot projections represent the same knot if and only if. up to planar isotopy, one can be obtained from the other via a sequence of \Index{Reidemeister moves}.
\end{lemma}

\todo[inline]{draw a picture}

% section intro_knots (end)


\section{Algebraic number fields and ring of integers} % (fold)
\label{sec:rings_and_ideals}
Throughout this section any ring is considered to be commutative with identity element; any ring homomorphism $\varphi \colon R \to S$ should satisfy $\varphi(1_R) =1_S$.

\begin{definition}[{name=[Integral Elements]{\cite[Def.~2.1]{neukirch}}}]
	Let $A \subseteq B$ be an extension of rings.
	An element $b \in B$ ist called \Index{integral} over $A$, if ist satisfies a monic equation
	\[
		x^n + a_1 x^{n-1} + \ldots + a_n =0
	\]
	with coefficients $a_i \in A$, $n\ge 1$.
	$B$ is called \Index{integral} over $A$, if all its elements are integral.
\end{definition}

\begin{remark}
	The sum and the product of two integral elements are again integral.
	While this might seem reasonable, the proof is far from being trivial! (see \textcite[Prop.~2.2]{neukirch})
\end{remark}

\begin{definition}[{name=[integral closure]}]
	The \emph{ring}
	\[
		\overline{A} \coloneqq \set{b \in B \given b \text{ integral over }A}
	\]
	is called \Index{integral closure} of $A$ in $B$.
	$A$ is called \Index{integrally closed}, if $\overline{A} = A$.
\end{definition}

By proving the transitivity of the relation \enquote{is integral over} one sees, that $\overline{A}$ is indeed integrally closed.

\begin{definition}[{name=[algebraic number field]}]
	A finite field extension of $\mathbb{Q}$ is called an \Index{algebraic number field}.
\end{definition}

This leads us to an object of central interest in algebraic number theory:

\begin{definition}[{name=[ring of integers]}]
	Let $K$ be some algebraic number field.
	The integral closure $\mathcal{O}_K \subseteq K$ of $\mathbb{Z} \subseteq \mathbb{Q}$ is called \Index{ring of integers} of the algebraic number field $K$.
\end{definition}

For fans of commutative diagrams:
\[
	\begin{tikzcd}
		\mathcal{O}_K \rar[hook] & K \\
		\mathbb{Z} \uar[hook] \rar[hook] & \mathbb{Q} \uar[hook]
	\end{tikzcd}
\]
We think of $\mathcal{O}_K$ as a generalization of $\mathbb{Z}$, which raises the question whether $\mathcal{O}_K$ has a \Index{unique prime factorization}. Let's consinder an example.
Let $D\neq 0$ be some square-free integer.
The field
\[
	K=\mathbb{Q}(\sqrt{D}) = \set*{a + b \sqrt{D} \given a,b \in \mathbb{Q}}
\]
is a \Index{quadratic field}. $\mathcal{O}_K$ can be explicitly computed:

\begin{lemma}\label{lem:quadratic_integers}
	Let $D\neq 0$ be some square-free integer.
	With $K=\mathbb{Q}(\sqrt{D})$ we have
	\[
		\mathcal{O}_K = \begin{cases}
			\mathbb{Z}\benbrace[\big]{\sqrt{D}}, &\text{ if }D \equiv 2,3 \mod 4\\[1ex]
			\mathbb{Z}\benbrace[\big]{\frac{1+ \sqrt{D}}{2}}, &\text{ if } D \equiv 1 \mod 4
		\end{cases}
	\]
\end{lemma}
\begin{proof}
	If $D \equiv 1 \mod 4$ then $p(X)= X^2 -X + (1-D)/4 \in \mathbb{Z}[X]$ and $p \enbrace[\big]{\frac{1+ \sqrt{D}}{2}}=0$, therefore $\frac{1+ \sqrt{D}}{2}$ is integral.
	So $\mathcal{O}_K$ has $\mathbb{Z}[\sqrt{D}]$ as a subring and if $D \equiv 1 \mod 4$ the bigger subring $\mathbb{Z}\benbrace[\big]{\frac{1+ \sqrt{D}}{2}}$.
	We show, that there are no other integral elements.
	The minimal polynom of $a + b \sqrt{D}$ with $b\neq 0$ is
	\[
		p(X) = X^2 -2aX + (a^2 - Db^2)
	\]
	So for $a +b \sqrt{D}$ to be integral we need $2a \in \mathbb{Z}$ and $(a^2 - Db^2) \in \mathbb{Z}$.
	Since $D$ is square-free we also get $2b \in \mathbb{Z}$ and $(a \in \mathbb{Z} \implies b \in \mathbb{Z})$.

	This leaves us with the $a \notin \mathbb{Z}$, implying that $2a=2k+1$ is odd.
	We get
	\[
		a^2 + Db^2 = \enbrace*{\frac{2k+1}{2}}^2 + \frac{4Db^2}{4} = \frac{4j^2 +4j + 4 Db^2 + 1}{4}
	\]
	The numerator has residue $1$ modulo $4$ while the denominator is divisible by $4$, if $b \in \mathbb{Z}$.
	Hence $b \notin \mathbb{Z}$ and $2b =2 l + 1$, which leads us to
	\[
		a^2 + Db^2 = \frac{4k^2 +4k +1 - D(4l^2 + 4l +1)}{4}
	\]
	For this to be an integer, we need to have $D \equiv 1 \mod 4$.
	Our element $a + b \sqrt{D}$ is contained in the ring $\mathbb{Z}\benbrace[\big]{\frac{1+ \sqrt{D}}{2}}$ as a simple application of the famous \enquote{add zero}-trick shows:
	\[
		\frac{2k+1}{2} + \frac{2l+1}{2} \cdot \sqrt{D} = \frac{2k+1}{2} - \frac{2l+1}{2} + \frac{2l+1}{2} + \frac{(2l+1)\sqrt{D}}{2} = \Underbrace{(k-l)}{\in \mathbb{Z}} + \Underbrace{(2l+1)}{\in \mathbb{Z}} \cdot \frac{1+ \sqrt{D}}{2}
	\]
	This finished the proof.
\end{proof}

\begin{example}
	For $d=-5$ we have $\mathcal{O}_K = \mathbb{Z}[\sqrt{-5}] = \mathbb{Z} + \mathbb{Z}\sqrt{-5}$ by \cref{lem:quadratic_integers}.
	The following factorization into irreducible elements is not unique
	\[
		21 = 3 \cdot 7 = \enbrace*{1+ 2 \sqrt{-5}} \cdot \enbrace*{1 - 2 \sqrt{-5}}
	\]
\end{example}

Historically this failure of of unique factorization led to the discovery of ideal theory (by \textsc{Edward Kummer} 1847).
Some basic properties of $\mathcal{O}_K$ concerning ideals are:

\begin{theorem}[{name={\cite[Thm.~3.1]{neukirch}}}]
	The ring $\mathcal{O}_K$ is noetherian, integrally closed and every prime ideal $\mathfrak{p} \neq 0$ is a maximal ideal.
\end{theorem}

This set of properties leads us to the following definition:\footnote{noetherian: every ascending chain of ideals is finite $\iff$ every ideal is finitely generated}

\begin{definition}[{name=[Dedekind domain]}]
	A noetherian, integrally closed integral domain in which every nonzero prime ideal is maximal is called a \Index{Dedekind domain}.
\end{definition}

Dedekind domains can be viewed as a generalization of principal ideal domains --- just as $\mathcal{O}_K$ may be viewed as a generalization of $\mathbb{Z}$.
The notion of divisibility for ideals is the following: Let $\mathfrak{a}, \mathfrak{b} \subseteq \mathcal{O}$ be ideals in a Dedekind domain.
Then $\mathfrak{a} \mid \mathfrak{b}$ is defined as $\mathfrak{b} \subseteq \mathfrak{a}$, the sum of ideals is
\[
	\mathfrak{a} + \mathfrak{b} = \set*{a + b \given a \in \mathfrak{a}, b \in \mathfrak{b}}
\]
and may be regarded as the greatest common divisor.
The least common multiple is the intersection $\mathfrak{a} \cap \mathfrak{b}$.
The product of two ideals is
\[
	\mathfrak{a} \mathfrak{b} = \set*{\sum\nolimits_i a_i b_i \given a_i \in \mathfrak{a}, b_i \in \mathfrak{b}}
\]
The ideals with this multiplication now have a unique prime factorization:

\begin{theorem}[{name={\cite[Thm.~3.3]{neukirch}}}]
	Every ideal $\mathfrak{a}$ of a Dedekind domain $\mathcal{O}$ different from $(0)$ and $(1)$ admits a factorization
	\[
		\mathfrak{a} = \mathfrak{p}_1 \cdots \mathfrak{p}_r
	\]
	into nonzero prime ideals $\mathfrak{p}_i$ of $\mathcal{O}$ which is unique up to the order of the factors.
\end{theorem}

There is a notion of dimension for rings which we will need later on
\begin{definition}[{name=[Krull dimension]}]
	The \Index{Krull dimension} of a commutative ring is defined as the supremum of integers $n$ such that there is a strict chain of prime ideals
	\[
		\mathfrak{p}_0 \subsetneq \mathfrak{p}_1 \subsetneq \mathfrak{p}_2 \subsetneq \cdots \subsetneq \mathfrak{p}_n
	\]
\end{definition}

Obviously every Dedekind domain has Krull dimension one.

\section{Localization}
\label{sec:localization}
Given a integral domain $R$ it is easy to construct the \Index{field of fractions}
\[
	K = \set*{\frac{r}{s} \given r \in R, s \in R \setminus \set*{0}}
\]
This construction can be generalized by choosing any nonempty $S \subseteq R \setminus \set*{0}$ which is closed under multiplication instead of $R \setminus \set*{0}$.
\[
	S^{-1}R = \set*{\frac{r}{s} \given r \in R, s \in S}
\]
This is called the \Index{localization} of $R$ at $S$ and has a canonical ring structure.
By the very definition of prime ideal the set $S=R \setminus \mathfrak{p}$ is closed under multiplication.
In this special case one usually writes $R_{\mathfrak{p}}$ instead of $S^{-1}R$.

The localization $R_{\mathfrak{p}}$ is often used when dealing with a single prime ideal $\mathfrak{p}$, because $R_{\mathfrak{p}}$ \textquote[\cite{neukirch}]{forgets everything that has nothing to do with $\mathfrak{p}$}, which is illustrated by the following lemmas:

\begin{lemma}
	The mappings
	\[
		\mathfrak{q} \longmapsto S^{-1}\mathfrak{q} \qquad \mathfrak{Q} \longmapsto \mathfrak{Q} \cap R
	\]
	are mutually inverse one-to-one correspondences between the prime ideals $\mathfrak{q} \subseteq R \setminus S$ of $A$ and the prime ideals $\mathfrak{Q}$ of $S^{-1}R$.
\end{lemma}

\begin{lemma}
	$R_{\mathfrak{p}}$ is a \Index{local ring}, i.e. $R_{\mathfrak{p}}$ has a unique maximal ideal, namely $\mathfrak{m}_{\mathfrak{p}} = \mathfrak{p} R_{\mathfrak{p}}$.
	There is a canonical embedding
	\[
		R/\mathfrak{p} \hookrightarrow R_{\mathfrak{p}}/\mathfrak{m}_{\mathfrak{p}}
	\]
	indentifying $R_{\mathfrak{p}}/\mathfrak{m}_{\mathfrak{p}}$ with the field of fractions of $R/\mathfrak{p}$.
\end{lemma}

\todo[inline]{maybe talk about discrete valuation rings}




\section{Affine schemes -- the geometrization of ring theory}
\label{sec:affine schemes}

Let's begin with a simple yet intructive example:
\begin{example}
	\label{ex:complex_poly}
	Let $f$ be some complex polynomial.
	$f$ may be interpreted as a function on the complex plane and for a point $a$ in the complex plane the set of all functions that vanish at $a$ is a maximal ideal $\mathfrak{p}=(X-a)$ of $\mathbb{C}[X]$.
	We get a one-to-one correspondence between the complex plane and the set $M$ of maximal ideals of $\mathbb{C}[X]$.
	An element $f \in \mathbb{C}[X]$ may now be viewed as an function on $M$ as follows:
	For every point $\mathfrak{p}=(X-a)$ of $M$ we have a canonical isomorphism
	\[
		\begin{tikzcd}
			\mathbb{C}[X]/\mathfrak{p} \rar["\sim"] & \mathbb{C}
		\end{tikzcd}
	\]
	induced by evaluation at $a$.
	We may thus view the residue class
	\[
		f(\mathfrak{p}) \coloneqq f \mod \mathfrak{p} \in \kappa(\mathfrak{p})
	\]
	in the residue class field $\kappa(\mathfrak{p}) = \mathbb{C}[X]/\mathfrak{p}$ as the \enquote{value} of $f$ at the point $\mathfrak{p} \in M$.
\end{example}

\begin{definition}[{name=[prime spectrum]}]
	Let $R$ be a commutative ring.
	\[
		\Spec(R) \coloneqq \set*{\mathfrak{p} \subset R \given \mathfrak{p} \text{ is a prime ideal} }
	\]
	is called the \Index{prime spectrum} of $R$.
	$\Spec(R)$ is endowed with the \Index{Zariski topology} defined by requiring the following sets to be closed
	\[
		V(\mathfrak{a}) = \set*{\mathfrak{p} \given \mathfrak{p} \supseteq \mathfrak{a}}\marginnote{note that $V(\mathfrak{a}) \cup V(\mathfrak{b}) = V(\mathfrak{a} \mathfrak{b})$}
	\]
	where $\mathfrak{a}$ varies over the ideals of $R$.
	A basis is given by the sets $U_f = \set*{\mathfrak{p} \given f \notin \mathfrak{p}}$.
\end{definition}

Note that $X= \Spec(R)$ is usually not Hausdorff and that the closed points correspond to the maximal ideals of $R$.
Pickung up \cref{ex:complex_poly} we now can define the \enquote{value} of $f \in R$  as
\[
	f(\mathfrak{p}) \coloneqq f \mod \mathfrak{p}
\]
in $\kappa(\mathfrak{p})$ the field of fractions of $R/\mathfrak{p}$.
Therefore the values of $f$ do not lie in the same field in general!
The zero ideal $\mathfrak{p}= (0)$ is called the \Index{generic point} of $X$, since its closure in the Zariski topology is the total space $X$.
The sets $V(\mathfrak{a})$ are the vanishing sets of ideals $\mathfrak{a}$.
Let's consider another example

\begin{example}
	$X= \Spec(\mathbb{Z})$ may be represented by a line, where for every prime number there is a closed point and there is the generic point $(0)$ as well.
	Since
	\[
		X \setminus V(\mathfrak{a}) = \set*{\mathfrak{p} \given \mathfrak{p} \not\supseteq \mathfrak{a}}
	\]
	an open set is obtained by by throwing out finitely many prime numbers $p_1, \ldots, p_n$.
	The value of $a \in \mathbb{Z}$ at the point $(p)$ is the residue class
	\[
		a(p) = a \mod p \in \mathbb{Z}/p \mathbb{Z}
	\]
	Therefore the field\emph{s} of of values are then $\mathbb{Z}/2 \mathbb{Z}, \mathbb{Z}/3 \mathbb{Z}, \ldots , \mathbb{Q}$.
\end{example}

This geometric interpretation of ring elements as functions on $X=\Spec(R)$ can be further refined by the so called \Index{structure sheaf} $\mathcal{O}_X$:
For simplicity we restrict our attention to one-dimensional integral domains $R$.
We associate to ever nonempty open set $U \subseteq X = \Spec(R)$ the ring of \enquote{regular functions} on $U$
\[
	\mathcal{O}(U) = \set*{\frac{f}{g} \given g(\mathfrak{p}) \neq 0 \text{ for all } \mathfrak{p} \in  U}
\]
This is the localization of $R$ with respect to the multiplicative set $S = R \setminus \bigcup_{\mathfrak{p} \in U} \mathfrak{p}$.
There are \Index{restriction maps} $\rho_{UV} \colon \mathcal{O}(U) \to \mathcal{O}(V)$ which turn the system of rings $\mathcal{O}(U)$ into a sheaf on $X$.

\begin{definition}[{name=[presheaf]}]
	A \Index{presheaf} $\mathcal{F}$ of rings is a contravariant functor
	\[
		\mathcal{F} \colon X_{\mathrm{top}} \longrightarrow \textsc{Rings}
	\]
	such that $\mathcal{F}(\emptyset)=0$, where $X_{\mathrm{top}}$ is the category which objects are the open sets of $X$ and the morphisms are given by inclusion.
	An element in $\mathcal{F}(U)$ is called \Index{section} over $U$.
\end{definition}

\begin{definition}[{name=[sheaf]}]
	A presheaf $\mathcal{F}$ on the topological space $X$ is called a \Index{sheaf}, if for all open coverings $\set*{U_i}$ of any open set $U$ one has
	\begin{itemize}[(i)]
		\item If $s,s' \in \mathcal{F}(U)$ with $s|_{U_i} = s'|_{U_i}$ for all $i \implies s=s'$.
		\item If $s_i \in \mathcal{F}(U_i)$ is a family of sections such that $s_i|_{U_i \cap U_j} = s_j|_{U_i \cap U_j}$ for alle $i,j$, then there is a section $s \in \mathcal{F}(U)$ such that $s|_{U_i} = s_i$ for all $i$.
	\end{itemize}
\end{definition}

The structure sheaf $\mathcal{O}(X)$ for $X=\Spec(R)$ is in fact a sheaf and the couple $(X, \mathcal{O}_X)$ is called an \Index{affine scheme}.

\subsection*{Quotients correspond to closed subschemes}
\label{sub:quotients_corres}
Let $I \subseteq R$ be some ideal.
The projection $s \colon R \twoheadrightarrow R/I$ induces a continous map
\mapdef{s^* \colon \Spec(R/I)}{\Spec{R}}{\mathfrak{p}}{s^{-1}(\mathfrak{p})}{}
It is a basic algebra excercise to prove, that this is well defined and the image consists exactly of the prime ideals, that contain the ideal $I$.
The image is therefore closed in the Zariski topology and gives rise to a closed subscheme of the affine scheme.

\begin{definition}[{name=[stalk and germs]}]
	The \Index{stalk} of the sheaf $\mathcal{F}$ at the point $x \in X$ is defined to be the direct limit 
	\[
		\mathcal{F}_x = \varinjlim_{U \ni x} \mathcal{F}(U)
	\]
	where $U$ varies over all open neighbourhoods of $x$.
	Let $s_U \in \mathcal{F}(U)$ and $s_V \in \mathcal{F}(V)$ be sections.
	They are equivalent if there is a neighbourhood $W \subseteq U \cap V$ of $x$ such that $s_U|_W = s_V|_W$.
	Elements of $\mathcal{F}_x$ are such equivalence classes and called \Index{germs} of sections at $x$.
\end{definition}

The correspondence between completions and \enquote{inifitesimal neighbourhoods} can not be made completely explicit right now, but we consider an example in the next section. 

\section{$p$-adic numbers}
\label{sec:p-adic numbers}

\begin{definition}
	Fix a prime number $p$.
	The \Index{$p$-adic integers} are defined as the projective limit 
	\[
		\mathbb{Z}_p = \varprojlim \mathbb{Z}/p^n \mathbb{Z}
	\]
	This is the \Index{completion} at the ideal $(p)$.
\end{definition}
This means a $p$-adic integer is a sequence $(a_n)_n$ of residue classes $a_n \in \mathbb{Z}/p^n \mathbb{Z}$ which are compatible in the following sense:
Let $\lambda_n \colon \mathbb{Z}/p^{n+1} \mathbb{Z} \twoheadrightarrow \mathbb{Z}/p^n \mathbb{Z}$ be the canonical projection.
Then $\lambda_n(a_{n+1}) = a_n$ for all $n \in \mathbb{N}$, i.e. the projective system looks like this
\[
	\begin{tikzcd}
		\mathbb{Z}/p \mathbb{Z} & \mathbb{Z}/p^2 \mathbb{Z} \lar["\lambda_1"'] & \mathbb{Z}/p^3 \mathbb{Z} \lar["\lambda_1"'] & \ldots \lar["\lambda_3"']
	\end{tikzcd}
\]
The projective limit is a subring of the product ring $\prod_{n \in \mathbb{N}} \mathbb{Z}/p^n \mathbb{Z}$, therefore addition and multiplication are defined componentwise.
The additive inverse to $(a_n)_n$ is $(p^n - a_n)_n$ and every sequence with $a_1 \neq 0$ has a multiplicative inverse, since $a_1 \neq 0$ implies, that every $a_n$ is not divisible by $p^n$ and therefore as an inverse $b_n$.
The sequence $(b_n)_n$ is an inverse for $(a_n)_n$.

\begin{definition}
	The $p$-adic numbers $\mathbb{Q}_p$ are defined as the field of fractions of $\mathbb{Z}_p$.
\end{definition}

\begin{itemize}
	\item The mapping $a \mapsto (a \mod p^n)_n$ defines an embedding $\mathbb{Z} \hookrightarrow \mathbb{Z}_p$.
	\item Every rational number $x$ can be written as $x = \pm \frac{a}{b} p^n$ for a unique $n$ and $p \nmid ab$.
	The \Index{$p$-adic absolute value} of $x$ is then given by
	\[
		\abs{x}_p \coloneqq p^{-n} \qquad \abs*{0}_p \coloneqq 0
	\]
	It can be easily shown, that this is a norm.
	\item The completion of $\mathbb{Q}$ with respect to this norm is again $\mathbb{Q}_p$, since the completion consists of equivalence classes of Cauchy sequences in $\mathbb{Q}$ and the $p$-adic value of $x$ is small, if $x$ is divisible by a high power of $p$.
	\item $\mathbb{Z}_p$ is an integral domain: If $ab=0$ we get $\abs*{a}_p \cdot \abs*{b}_p = 0$ and therfore $a=0$ or $b=0$.
	\item $\mathbb{Z}_p$ is a local ring (even a discrete valuation ring) with maximal ideal $\mathfrak{m} = p \mathbb{Z}_p$.
	The residue class field is $\mathbb{Z}_p/p \mathbb{Z}_p \cong \mathbb{Z}/p \mathbb{Z} = \mathbb{F}_p$ (Iso with isomorphism theorem).
	On the spectras the inclusion $\mathbb{Z} \hookrightarrow \mathbb{Z}_p$ induces a map $\Spec(\mathbb{Z}_p) \to \Spec(\mathbb{Z})$, which extracts the scheme $\Spec(\mathbb{Z}_p)$ from the line $\Spec(\mathbb{Z})$.
	\todo[inline]{draw picture}
\end{itemize}



% section rings_and_ideals (end)



\cleardoubleoddemptypage%
\pagenumbering{Alph}
\setcounter{page}{1}
\appendix%
\printbibliography%
\printindex%
\todototoc%
\listoftodos[To-do's]
\end{document}
